\documentclass[11pt]{article}
\usepackage[T1]{fontenc}
\usepackage{xcolor}
\usepackage{amsmath}
\usepackage[colorlinks=true,linkcolor=black,citecolor=black]{hyperref}

% Revised text is blue in the review version. Set it to black for a clean build;
% NotePDF finds the text either way.
\definecolor{revcolor}{RGB}{0,0,200}
\newcommand{\rev}[1]{{\color{revcolor}#1}}   % detected automatically as a marker
\newcommand{\ours}{FlowNote\xspace}
\providecommand{\xspace}{}

\title{FlowNote: Annotating Revisions Automatically}
\author{A. Author}

\begin{document}
\maketitle

\section{Introduction}
Reviewers often ask authors to point out what changed between two versions of a paper.
{\color{revcolor} %@note: R1.1 -- clarified the motivation
We propose \ours, a tool that turns the revision markers already present in the \LaTeX{} source
into highlights and comments in the compiled PDF, so that reviewers can find every change
without reading the whole paper again~\cite{knuth,lamport}.}
Manual highlighting is slow and error-prone.

%@note: R2.3 -- new evaluation
\rev{We further evaluate \ours on three new datasets and report a speed-up of $3.2\times$ over manual annotation.}

\section{Method}
\ours works in three steps.
\begin{itemize}
  \item It reads the source and follows every \verb|\input| command.
  \item {\color{revcolor} %@note: R3.2 -- explain the matching step
  It converts each marked region to plain text and searches for it in the PDF, ignoring
  punctuation, hyphenation and line breaks.}
  \item It writes standard PDF annotations that every viewer can display.
\end{itemize}

{\color{revcolor}
The matching score is defined as
\begin{equation}
  s = \frac{|m|}{|t|},
\end{equation}
where $m$ is the matched text and $t$ the target text (see Section~\ref{sec:method}).}
\label{sec:method}


\section{Conclusion}
\textcolor{revcolor}{%@note: R1.4 -- added limitations
A limitation is that display equations are not highlighted automatically\footnote{They are listed in the report instead.}; we plan to address this in future work.}

\begin{thebibliography}{9}
\bibitem{knuth} D.~E. Knuth. \emph{The \TeX book}. Addison-Wesley, 1984.
\bibitem{lamport} L.~Lamport. \emph{\LaTeX: A Document Preparation System}. Addison-Wesley, 1994.
\end{thebibliography}
\end{document}
